\subsection{Envelope-mean subharmonicity}\label{supp:envelope-proof}
The envelope mean is
\[
\overline{\mathcal A}(\mathbf x)=\frac{1}{2\pi}\int_0^{2\pi}
\sqrt{\varepsilon_1(\mathbf x)^2+\varepsilon_2(\mathbf x)^2
+2\varepsilon_1(\mathbf x)\varepsilon_2(\mathbf x)\cos\vartheta}\,\mathrm d\vartheta,
\]
where $\vartheta$ is beat phase.

For fixed scalar coupling in homogeneous, isotropic, source-free tissue,
the envelope mean is subharmonic. At each beat phase $\vartheta$,
$\mathcal A$ is the Euclidean norm of
$(\varepsilon_1+\varepsilon_2\cos\vartheta,\,\varepsilon_2\sin\vartheta)$,
a vector of harmonic functions. A norm is convex, so this envelope is subharmonic;
averaging over $\vartheta$ preserves that property. Thus both idealized mean maps
lack a strict interior maximum under these conditions (Fig.~\ref{fig:spatial_maps}b).

\paragraph{Interior maxima of the excursion and the beat amplitude.}\label{supp:interior-maxima}
The product of the polarizations obeys
$\nabla^2(\varepsilon_1\varepsilon_2)=2\nabla\varepsilon_1\cdot\nabla\varepsilon_2$,
whose sign is unrestricted, and an excursion, a temporal maximum minus a temporal
minimum, need not be subharmonic. A local example shows that both can have a strict
three-dimensional interior maximum under the conditions above. In dimensionless
coordinates $\mathbf x=(x,y,z)$, take $0<c<1/2$, $q=x^2-\tfrac12(y^2+z^2)$, and the harmonic potentials
\[
\phi_{1,2}=-x\mp\tfrac12(x^2-y^2)-c\left[\tfrac13x^3-\tfrac12x(y^2+z^2)\right].
\]
With unit coupling along the fixed $x$ direction,
$\varepsilon_{1,2}=-\partial_x\phi_{1,2}=1\pm x+cq$, positive near the origin. Then
\[
\varepsilon_1\varepsilon_2=1+(2c-1)x^2-c(y^2+z^2)+O(|\mathbf x|^4),
\]
with a negative-definite Hessian at the origin, and the envelope excursion
$2\min(\varepsilon_1,\varepsilon_2)=2(1+cq-|x|)$ is also strictly maximal there.
This is an existence example, not a prediction for a montage. The subharmonic maps
above can still have interior minima; the restriction concerns maxima only.
